
\subsubsection{具体分析}

问题三要求在导通概率不低于90\%的前提下确定金属A的最小填充率，属于在蒙特卡洛模拟框架上的优化搜索问题。本问采用二分搜索策略在颗粒数空间中定位阈值。

基于问题二的结果，导通概率在0.50\%时已达100\%，而在0.0693\%（附件2）时为0\%，因此渗透阈值位于两者之间。首先通过粗测试确定合理的搜索区间，接着采用二分搜索在颗粒数空间$[5, 80]$中迭代定位使导通概率恰好达到90\%的临界颗粒数，然后通过大样本蒙特卡洛验证确保结果的统计可靠性。具体分析流程如图\ref{fig:analysis3_flow}所示。

\begin{figure}[ht]
  \begin{center}
    \begin{tikzpicture}[x=1cm, y=0.7cm]
      \node[box] (n11) at (0, 0) {确定区间};
      \node[box] (n12) at (5, 0) {二分搜索};
      \node[box] (n13) at (10, 0) {MC估概率};
      \node[box] (n22) at (5, -3) {区间更新};
      \node[box] (n21) at (0, -3) {收敛判断};
      \node[box] (n23) at (10, -3) {大样验证};
      \node[box] (n31) at (0, -6) {概率曲线};
      \node[box] (n32) at (5, -6) {临界分析};
      \node[box] (n33) at (10, -6) {结果输出};
      \draw[arrow] (n11) -- (n12);
      \draw[arrow] (n12) -- (n13);
      \draw[arrow] (n23) -- (n22);
      \draw[arrow] (n22) -- (n21);
      \draw[arrow] (n13) -- ++(0,-1.0) -| (n23);
      \draw[arrow] (n21) -- (n31);
      \draw[arrow] (n31) -- (n32);
      \draw[arrow] (n32) -- (n33);
    \end{tikzpicture}
    \caption{问题三分析流程图}
    \label{fig:analysis3_flow}
  \end{center}
\end{figure}

\subsubsection{模型准备}

在建立搜索模型之前，需要确定搜索空间和分析附件3的参考配置。

根据问题二的初步测试，10个颗粒时导通概率约44\%，50个颗粒时导通概率约96\%。因此将二分搜索的初始区间设为颗粒数$[5, 80]$，对应的填充率区间约为$[0.0071\%, 0.1131\%]$。搜索精度设为$10^{-4}$（填充率0.0001\%），二分搜索约需$\log_2(75/0.0001) \approx 20$轮迭代。

附件3包含535个金属A颗粒的坐标数据，填充率为$535 \times 1.4137 \times 10^7 / 10^{12} = 0.7563\%$，远高于预期的渗透阈值，预期该配置应为导通状态。

该问题本质上是在随机模拟框架下的统计优化问题，每次目标函数评估（导通概率估计）都涉及蒙特卡洛模拟的随机误差，需要在搜索效率和统计精度之间取得平衡。综上所述，本节完成了搜索空间的界定，接下来将建立二分搜索模型并进行求解。
